\section{Extended related work}
\label{app:references}

This appendix collects the literature the benchmark rests on, grouped by the role each work plays
rather than alphabetically. Section~\ref{sec:related} discusses the works the design depends on
directly; the rest are recorded here.

\paragraph{Origami benchmarks and systems for language models.}
\citet{gamibench} releases both code and data, and evaluates viewpoint consistency and the
recognition of infeasible patterns. \citet{origamispace} evaluates pattern prediction, multi-step
spatial reasoning and crease-pattern code generation, and ships an interactive environment for the
last of these. \citet{origamibench} synthesises flat-foldable patterns in an interactive
environment; its instances are the 366
% fact:corpus.instagram.total
patterns distributed with Flat-Folder, and it is one source of the metrics surveyed while our
scoring was designed. \citet{corigami} co-designs flat-foldable and recognisable models, scoring
candidates by flat-foldability as a boolean and by a vision-language aesthetic judgement.
\citet{learn2fold} plans structured generation against a learned world model.
\citet{foldingagent} reconstructs folding processes from instructional video, over the
sequence-level corpus of \citet{purelandfold}.

\paragraph{Complexity and theory of flat folding.}
\citet{bern-hayes} carries the claim the task rests on: geometry does not determine layer order,
and deciding the overlap order of a flat folding is NP hard even when a valid mountain-valley
assignment is given. \citet{arkin-map} defines simple foldability and the one-layer, some-layers
and all-layers models this action space is drawn from, and \citet{akitaya-hard} settles the
complexity of deciding it for square paper with creases at multiples of $45^\circ$.
\citet{akitaya-infinite} treats the all-layers model, which is also the model of sheet-metal
bending, and \citet{mixed-orthogonal} treats mixed orthogonal crease patterns.
\citet{akitaya-flatfolder} decides whether a pattern admits a flat-folded state at all.
\citet{layer-algebra} gives an algebraic treatment of the layer-ordering constraints, and
\citet{flat-folding-graphs} studies the graphs of flat foldings themselves.
\citet{continuous} separates a folded state from a folding motion, with \citet{reaching-folded}
as its predecessor: reachability is free, so the hard question is the discrete step structure.
\citet{turing} establishes that flat origami is Turing complete, which bounds what any general
claim about foldability can promise.

\paragraph{The crease-pattern-to-sequence problem itself.}
\citet{akitaya-cp2seq} named the problem and supplied the tools the area still uses: reflection
paths, graph rewriting and step graphs. Its reported costs make the case for a benchmark as well
as any argument here does, with a frog base expanding to 22,665 nodes and half an hour of search,
and its future work asks for exactly the priority heuristic the models in this paper are being
asked to supply. \citet{creasy} is an open-source implementation of that line and is usable as a
symbolic baseline; Section~\ref{sec:exp-baselines} lists it among the baselines not yet run.

\paragraph{Counting folded states.}
The number of ways to fold a strip of $n$ labelled stamps runs 1, 2, 6, 16, 50, 144, 462, 1392,
4536, 14060 and has no closed form \citep{oeis-a000136}. The problem is old \citep{lucas} and has
been attacked repeatedly since \citep{koehler,lunnon}; meanders and stamp foldings are the same
combinatorial object \citep{meanders}. The relevance here is that the number of valid states grows
fast enough that scoring against any single reference sequence would be wrong, which is why
Section~\ref{sec:scoring} scores by execution instead.

\paragraph{Spatial reasoning context.}
\citet{spatial-survey} surveys spatial reasoning in multimodal models. \citet{vot} names the
mental-imagery question this benchmark turns into a measurement: whether a model can maintain and
update a spatial representation across steps rather than describe one.

\paragraph{Statistical physics of random flat foldability.}
That one crease pattern admits many valid folded states is the normal case rather than a
curiosity, and the statistical-physics literature on random flat foldability reaches the same
conclusion from another direction \citep{spin-model,random-locally}.

\paragraph{Tools, formats and data.}
\citet{flatfolder} is the implementation that decides and enumerates flat-folded states, for the
method of \citet{akitaya-flatfolder}; together they are the source of the four constraint classes
cited in Section~\ref{sec:env-refusals}. \citet{fold-format} is the de facto standard used by
every artifact we ship. \citet{oripa} is a crease-pattern editor with folded-form estimation.
\citet{purelandfold} is the closest existing sequence-level corpus and is FoldingAgent's data; it
is not used here as training or evaluation data.

\paragraph{Scope, and what is left out.}
Thick folding, bar-and-hinge mechanics and the creased-sheet mechanics literature belong to a
different line of work, one in which the sheet has thickness and the verifier is a simulation
rather than a decision. That line asks a question this benchmark does not: the states here are
zero-thickness and flat, and the environment rules on a fold in one pass rather than integrating a
deformation. Those works are therefore not treated as neighbouring work, and the limits of the
model used here are stated in Section~\ref{sec:limitations}.
